\documentclass[10pt,a4paper]{amsart}
\usepackage[margin=19mm]{geometry}
\usepackage{amsmath,amssymb,mathtools}
\usepackage[scr=rsfs,cal=euler]{mathalfa}
\usepackage{xfrac}
\usepackage[unicode,hidelinks,bookmarks=false]{hyperref}
\newcommand{\ie}{i.e.\ }
\newcommand{\cf}{cf.\ }
\newcommand{\resp}{respectively\ }
\newcommand{\Subsection}{\subsection}
\providecommand{\nobreakdash}{\nobreak\hbox{-}}
\setlength{\emergencystretch}{2em}
\multlinegap0pt
\usepackage[shortlabels]{enumitem}
\usepackage{xcolor}
\usepackage{mathtools}
\usepackage{dsfont}
\newtheorem{lemma}{Lemma}
\newtheorem{proposition}{Proposition}
\newcommand{\Ex}{\mathbb{E}}
\title{On the intersection of pairs of trees}
\author{Mikl\'os B\'ona, Fabian Burghart, Stephan Wagner}
\date{Source: arXiv:2501.18570v4 (CC BY 4.0)}
\begin{document}
\maketitle

\noindent\emph{Source and licence.} On the intersection of pairs of trees. Version arXiv:2501.18570v4 (CC BY 4.0), distributed by arXiv under the Creative Commons Attribution 4.0 International licence, http://creativecommons.org/licenses/by/4.0/, as recorded in the arXiv licence metadata for this exact version and shown on the abstract page https://arxiv.org/abs/2501.18570v4. Full licence notice: \texttt{licenses/arxiv-2501.18570-CC-BY-4.0.txt}.\par\smallskip

\noindent\textit{Editorial scope.} Counted unit: the article's proposition prop:moments (source0.tex lines 126--130). The article's complete original proof of it is delivered with it (source0.tex lines 132--169, one \texttt{proof} environment of 38 lines). No mathematics is written, repaired or re-derived: the statement and the proof are the article's own lines, in the article's own notation, and no retained line is substituted or re-flowed.  Retained uncounted as the source-local context the proof needs: lines 111--113; lines 121--123.  Nothing was omitted from any retained span, so the package carries no omission record.  No retained line contains a citation, so the wrapper carries no bibliography and no bibliography entry.  No retained line contains a figure environment, an image-inclusion command or drawing code, so no image is delivered and the package declares no asset dependency.  Editorial formatting only, in addition: an AMS \texttt{amsart} preamble that restates the article's own declarations and macros used by the retained lines, namely the environments \texttt{lemma}, \texttt{proposition} on separate counters, together with the standard packages those lines need.  The retained environments print in this document as Lemma 1, Proposition 1 in order of appearance, whereas the article prints the same items with the numbering of its own full document.  The complete native source of the article as distributed in the arXiv e-print for this exact version is bundled unchanged as \texttt{sources/172467/source0.tex}.
\begin{lemma}\label{lem:given_components}
For a given forest $F$ on $[n]$ with components of sizes $n_1,n_2,\ldots,n_k$, the number of spanning trees containing $F$ as a subgraph is $(n_1n_2 \cdots n_k) n^{k-2}$.
\end{lemma}

\begin{equation}\label{eq:sum_of_ind}
X_n = \sum_e Z_e.
\end{equation}

\begin{proposition}\label{prop:moments}
For every fixed positive integer $r$, the $r$-th moment of $X_n$ is asymptotically equal to
\[e^{-2} \sum_{k \geq 0} \frac{2^k k^r}{k!} + O(n^{-1}),\]
i.e., it converges to the $r$-th moment of a Poisson random variable with expected value $2$.
\end{proposition}

\begin{proof}
By~\eqref{eq:sum_of_ind},
\[\Ex(X_n^r) = \sum_{e_1,e_2,\ldots,e_r}  \Ex(Z_{e_1}Z_{e_2} \cdots Z_{e_r}),\]
the sum being over all ordered $r$-tuples of (not necessarily distinct) edges. The summands can be simplified to terms of the form
\[\Ex(Z_{e_1}^{a_1} Z_{e_2}^{a_2} \cdots Z_{e_s}^{a_s}),\]
where the $e_i$ are distinct. Since the $Z_{e_i}$ are indicator random variables, this is equal to
\begin{equation}\label{eq:red_term}
\Ex(Z_{e_1}Z_{e_2}\cdots Z_{e_s}).
\end{equation}
The sum of the coefficients of all terms that reduce to~\eqref{eq:red_term} is
\[\sum_{\substack{a_1,a_2,\ldots,a_s \geq 1 \\ a_1 + a_2 + \cdots + a_s = r}} \binom{r}{a_1,a_2,\ldots,a_r}
= r! [t^r] (e^t-1)^s.\]
If $e_1,e_2,\ldots,e_s$ are disjoint edges, then by Lemma~\ref{lem:given_components}, we have
\[\Ex(Z_{e_1}Z_{e_2}\cdots Z_{e_s}) = \Big(\frac{2^s n^{n-s-2}}{n^{n-2}}\Big)^2 = 2^{2s} n^{-2s}.\]
Now consider an arbitrary set of $s$ distinct edges $e_1,e_2,\ldots,e_s$. If some of these edges form a cycle, then clearly $\Ex(Z_{e_1}Z_{e_2}\cdots Z_{e_s}) = 0$ as spanning trees cannot contain cycles. Otherwise, they induce a forest with $n-s$ components on the vertex set $[n]$. If $m_1,m_2,\ldots,m_{n-s}$ are its component sizes, then by Lemma~\ref{lem:given_components} there are
\[(m_1m_2\cdots m_{n-s})n^{n-s-2}\]
spanning trees that contain all the edges $e_1,e_2,\ldots,e_s$. By the elementary inequality $m \leq 2^{m-1}$ that is valid for all positive integers $m$, we have
\[(m_1m_2\cdots m_{n-s})n^{n-s-2} \leq 2^{m_1+m_2+\cdots+m_{n-s}-n+s} n^{n-s-2} = 2^s n^{n-s-2}.\]
It follows that
\[\Ex(Z_{e_1}Z_{e_2}\cdots Z_{e_s}) = 
\Big(\frac{(m_1m_2\cdots m_{n-s})n^{n-s-2}}{n^{n-2}}\Big)^2
\leq \Big(\frac{2^s n^{n-s-2}}{n^{n-2}}\Big)^2 = 2^{2s} n^{-2s},\]
thus $\Ex(Z_{e_1}Z_{e_2}\cdots Z_{e_s}) = O(n^{-2s})$ in all cases.
Of the $\binom{n(n-1)/2}{s}$ different sets of $s$ edges, the number of sets of disjoint edges is
\[\binom{n}{2s} (2s-1)!! = \frac{n(n-1)\cdots (n-2s+1)}{2^s s!},\]
since there are $\binom{n}{2s}$ choices for the vertices and $(2s-1)!!$ ways to pair them. Note that $\binom{n(n-1)/2}{s}$ and $\binom{n}{2s} (2s-1)!!$ are both $\frac{n^{2s}}{2^s s!} + O(n^{2s-1})$, implying that the number of sets of $s$ edges that are not all disjoint is $O(n^{2s-1})$. So it follows that
\begin{align*}
\Ex(X_n^r) &= \sum_{s=1}^r \sum_{\{e_1,e_2,\ldots,e_s\}} r! [t^r] (e^t-1)^s \Ex(Z_{e_1}Z_{e_2}\cdots Z_{e_s}) \\
&= \sum_{s=1}^r r! [t^r] (e^t-1)^s \Big( \binom{n}{2s} (2s-1)!! \frac{2^{2s}}{n^{2s}} + O \Big( n^{2s-1} \cdot n^{-2s} \Big) \Big) \\
&= \sum_{s=1}^r r! [t^r] (e^t-1)^s \Big( \frac{2^s}{s!} + O(n^{-1}) \Big).
\end{align*}
Thus the $r$-th moment of $X_n$ converges to
\begin{align*}
\sum_{s=1}^r r! [t^r] (e^t-1)^s \frac{2^s}{s!} &= r! [t^r] \sum_{s = 0}^{\infty} \frac{2^s (e^t-1)^s}{s!} = r! [t^r] e^{2(e^t-1)} \\
&= r! e^{-2} [t^r] \sum_{k \geq 0} \frac{(2e^t)^k}{k!} = e^{-2} \sum_{k \geq 0} \frac{2^k k^r}{k!},
\end{align*}
completing the proof.
\end{proof}

\end{document}
